\subsection{Normative Material and Specificity}

Architecture prose, ordered steps, architecture tables, instruction encoding diagrams, and the Operation
and Detailed Semantics (:term:base.terms.field|plural:) of instruction entries are normative unless a section explicitly identifies material
as a summary, example, or validation status.

\noindent
\begin{tabularx}{\linewidth}{@{}>{\raggedright\arraybackslash}p{1.55in}X@{}}
\toprule
\textbf{Rule set or provider} & \textbf{Covered material}\\
\midrule
Instruction definitions & Concrete instruction encodings, operand and effective-address grammar, extension wiring,
and document ordering.\\
\hyperref[page:instruction-execution-model]{Instruction Execution Model} common rules and instruction-entry Operation
and Detailed Semantics & Executable instruction behavior, (:term:base.terms.architectural_state:) transitions, fault and commit ordering,
repeat and architectural-event behavior, and single-logical-processor memory-access sequencing.\\
Formal memory model & Permitted cross-logical-processor ordering and visibility.\\
ABI specifications & ELF, C ABI, and calling-convention contracts.\\
Compiler-interface specifications & Source-language and compiler-facing target-interface contracts.\\
External numerical floating-point providers & Numerical results and certificates only; input validation and the resulting
architectural trap or commit behavior follow the applicable
\hyperref[page:instruction-execution-model]{Instruction Execution Model} common rules and instruction entry\textquotesingle{}s Operation
and Detailed Semantics.\\
\bottomrule
\end{tabularx}

Presented material remains normative where declared and follows the applicable
rule set. For executable behavior, readers apply the
\hyperref[page:instruction-execution-model]{Instruction Execution Model} common
rules together with the applicable instruction entry\textquotesingle{}s Operation and Detailed
Semantics. The following specificity rules govern narrowing among applicable
rules:

\begin{enumerate}
\item Opcode-space length, (:term:base.terms.required_instruction_length:), (:term:base.terms.encoded_instruction_length:), (:term:base.terms.field:) values, and decoding validity follow the selected instruction encoding, its
(:term:base.terms.field:) constraints, and the Instruction Encoding and Effective Addressing chapters.
\item Execution of a decoded instruction follows its Operation and Detailed Semantics. An instruction-specific rule
narrows the common execution, operand, flag, memory, or event rule that it names.
\item Common architectural behavior follows the chapter that defines that behavior. A structure-specific rule
narrows the defaults in Terminology and Compatibility.
\item Summary tables and examples provide navigation or explanatory context
and do not replace a normative rule.
\end{enumerate}

Two normative statements at the same specificity are required to agree. A discrepancy at that level is a
specification defect and is resolved by correcting the normative sources and their derived tables together.
